\textbf{Phần 1. Trắc nghiệm lựa chọn}
\setcounterpageref{ex}{0}

\begin{ex} %[3D2V58]
Tính giá trị của biểu thức: $A = 1250 \times 4 + 4800 : 6$.
\choice
{$5700$}
{\True $5800$}
{$5900$}
{$6000$}
\loigiai{
$1250 \times 4 = 5000$; $4800 : 6 = 800$; $5000 + 800 = 5800$.
\textbf{Đáp án của câu hỏi là B}
}
\end{ex}

\begin{ex} %[3D2V58]
Một cửa hàng có $7500\text{ kg}$ gạo. Ngày đầu bán được $1850\text{ kg}$, ngày thứ hai bán được gấp đôi ngày đầu. Cửa hàng còn lại bao nhiêu ki-lô-gam gạo?
\choice
{$3700\text{ kg}$}
{$2050\text{ kg}$}
{$1850\text{ kg}$}
{\True $1950\text{ kg}$}
\loigiai{
Ngày hai bán: $1850 \times 2 = 3700\text{ kg}$. Tổng hai ngày: $1850 + 3700 = 5550\text{ kg}$. Còn lại: $7500 - 5550 = 1950\text{ (kg)}$.
\textbf{Đáp án của câu hỏi là D}
}
\end{ex}

\begin{ex} %[3D2V58]
Tìm số tự nhiên $x$, biết: $(x - 1200) : 4 = 1500$.
\choice
{\True $7200$}
{$6800$}
{$7000$}
{$7400$}
\loigiai{
$x - 1200 = 1500 \times 4 = 6000 \Rightarrow x = 6000 + 1200 = 7200$.
\textbf{Đáp án của câu hỏi là A}
}
\end{ex}

\begin{ex} %[3D2V58]
Một mảnh vườn hình chữ nhật có chiều dài $240\text{ m}$, chiều rộng bằng $\dfrac{1}{3}$ chiều dài. Chu vi của mảnh vườn đó là
\choice
{$480\text{ m}$}
{$720\text{ m}$}
{\True $640\text{ m}$}
{$560\text{ m}$}
\loigiai{
Chiều rộng: $240 : 3 = 80\text{ m}$. Chu vi: $(240 + 80) \times 2 = 640\text{ (m)}$.
\textbf{Đáp án của câu hỏi là C}
}
\end{ex}

\begin{ex} %[3D2V58]
Tính nhanh giá trị: $2350 + 1800 + 1650 + 2200$.
\choice
{$7800$}
{\True $8000$}
{$8200$}
{$7900$}
\loigiai{
$(2350 + 1650) + (1800 + 2200) = 4000 + 4000 = 8000$.
\textbf{Đáp án của câu hỏi là B}
}
\end{ex}

\begin{ex} %[3D2V58]
Trong phép chia cho 7 có thương bằng 1250 và số dư lớn nhất có thể có. Số bị chia là
\choice
{$8750$}
{$8755$}
{$8757$}
{\True $8756$}
\loigiai{
Số dư lớn nhất là 6. Số bị chia: $7 \times 1250 + 6 = 8750 + 6 = 8756$.
\textbf{Đáp án của câu hỏi là D}
}
\end{ex}

\begin{ex} %[3D2V58]
Một xưởng dệt được $8400\text{ m}$ vải trong một tuần 6 ngày làm việc. Hỏi trong 4 ngày xưởng đó dệt được bao nhiêu mét vải?
\choice
{\True $5600\text{ m}$}
{$5400\text{ m}$}
{$5800\text{ m}$}
{$4200\text{ m}$}
\loigiai{
Mỗi ngày dệt: $8400 : 6 = 1400\text{ m}$. Bốn ngày dệt: $1400 \times 4 = 5600\text{ (m)}$.
\textbf{Đáp án của câu hỏi là A}
}
\end{ex}

\begin{ex} %[3D2V58]
Tổng của hai số là 4800, số lớn gấp 3 lần số bé. Số lớn là
\choice
{$1200$}
{$2400$}
{\True $3600$}
{$3200$}
\loigiai{
Tổng số phần: $3 + 1 = 4$ phần. Số bé: $4800 : 4 = 1200$. Số lớn: $1200 \times 3 = 3600$.
\textbf{Đáp án của câu hỏi là C}
}
\end{ex}

\begin{ex} %[3D2V58]
Một hình vuông có chu vi $4800\text{ cm}$. Diện tích của hình vuông đó là
\choice
{$1200\text{ cm}^2$}
{\True $1\,440\,000\text{ cm}^2$}
{$2400\text{ cm}^2$}
{$4800\text{ cm}^2$}
\loigiai{
Cạnh hình vuông: $4800 : 4 = 1200\text{ cm}$. Diện tích: $1200 \times 1200 = 1\,440\,000\text{ cm}^2$ (hoặc xét cạnh $12\text{ m}$).
\textbf{Đáp án của câu hỏi là B}
}
\end{ex}

\begin{ex} %[3D2V58]
Tính giá trị biểu thức: $(6500 - 2500) : 5 + 1400$.
\choice
{$2000$}
{$2100$}
{$2400$}
{\True $2200$}
\loigiai{
$4000 : 5 + 1400 = 800 + 1400 = 2200$.
\textbf{Đáp án của câu hỏi là D}
}
\end{ex}

\begin{ex} %[3D2V58]
Bác Nam mua một chiếc xe đạp giá $1\,850\,000$ đồng và đưa cho người bán hàng 2 triệu đồng. Bác Nam nhận lại số tiền thừa là
\choice
{\True $150\,000$ đồng}
{$250\,000$ đồng}
{$100\,000$ đồng}
{$200\,000$ đồng}
\loigiai{
Đổi 2 triệu đồng = $2\,000\,000$ đồng. Tiền thừa: $2\,000\,000 - 1\,850\,000 = 150\,000$ đồng.
\textbf{Đáp án của câu hỏi là A}
}
\end{ex}

\begin{ex} %[3D2V58]
Một đoàn gồm 5 xe tải nhỏ, mỗi xe chở $1450\text{ kg}$ hàng và 2 xe tải lớn, mỗi xe chở $2500\text{ kg}$ hàng. Cả đoàn xe chở tất cả bao nhiêu ki-lô-gam hàng?
\choice
{$11\,250\text{ kg}$}
{$12\,500\text{ kg}$}
{\True $12\,250\text{ kg}$}
{$10\,750\text{ kg}$}
\loigiai{
Xe nhỏ: $1450 \times 5 = 7250\text{ kg}$. Xe lớn: $2500 \times 2 = 5000\text{ kg}$. Tổng: $7250 + 5000 = 12\,250\text{ (kg)}$.
\textbf{Đáp án của câu hỏi là C}
}
\end{ex}

\textbf{Phần 2. Trắc nghiệm đúng sai}
\setcounterpageref{ex}{0}

\begin{ex} %[3D2V58]
Xét tính đúng sai của các bài toán tổng hợp phạm vi 10 000:
\choiceTF
{\True Biểu thức có ngoặc luôn được tính toán trong ngoặc trước.}
{$2000 \times 4 + 2000 = 2000 \times 4$.}
{\True $6400 : 8 \times 2 = 1600$.}
{Một số chia cho 8 dư 7 thì cộng thêm 1 đơn vị sẽ chia hết cho 8.}
\loigiai{
\begin{itemize}
    \item Ý a) Đúng: Quy tắc cơ bản.
    \item Ý b) Sai: Sai vì bằng $2000 \times 5$.
    \item Ý c) Đúng: $800 \times 2 = 1600$.
    \item Ý d) Sai: Mệnh đề đúng (thay mệnh đề sai): Một số chia cho 8 dư 7 thì cộng thêm 2 đơn vị chia hết cho 8.
\end{itemize}
\begin{itemize}
    \item a) \textbf{Đúng}.
    \item b) \textbf{Sai}.
    \item c) \textbf{Đúng}.
    \item d) \textbf{Sai}.
\end{itemize}
}
\end{ex}

\begin{ex} %[3D2V58]
Một nông trại nuôi 1500 con gà và 2500 con vịt:
\choiceTF
{\True Tổng số gà và vịt là 4000 con.}
{\True Số vịt nhiều hơn số gà là 1000 con.}
{Số vịt gấp đôi số gà.}
{\True Nếu bán đi 500 con gà thì số gà còn lại là 1000 con.}
\loigiai{
\begin{itemize}
    \item Ý a) Đúng: $1500 + 2500 = 4000$ con.
    \item Ý b) Đúng: $2500 - 1500 = 1000$ con.
    \item Ý c) Sai: Sai vì $1500 \times 2 = 3000 \neq 2500$.
    \item Ý d) Đúng: $1500 - 500 = 1000$ con.
\end{itemize}
\begin{itemize}
    \item a) \textbf{Đúng}.
    \item b) \textbf{Đúng}.
    \item c) \textbf{Sai}.
    \item d) \textbf{Đúng}.
\end{itemize}
}
\end{ex}

\textbf{Phần 3. Trả lời ngắn}
\setcounterpageref{ex}{0}

\begin{ex} %[3D2V58]
Tính giá trị của biểu thức: $1400 \times 4 - 3600 : 6$.
\shortans{5000}
\loigiai{
\textbf{Đáp số:} $5000$\\
$5600 - 600 = 5000$.
}
\end{ex}

\begin{ex} %[3D2V58]
Một hình chữ nhật có chiều dài $150\text{ m}$, chiều rộng $80\text{ m}$. Tính chu vi mảnh đất đó theo mét.
\shortans{460}
\loigiai{
\textbf{Đáp số:} $460$\\
$(150 + 80) \times 2 = 460\text{ (m)}$.
}
\end{ex}

\begin{ex} %[3D2V58]
Tính nhanh: $3250 + 1750 + 2400 + 1600$.
\shortans{9000}
\loigiai{
\textbf{Đáp số:} $9000$\\
$5000 + 4000 = 9000$.
}
\end{ex}

\begin{ex} %[3D2V58]
Một xưởng may có 4800 mét vải. Người ta đã may hết $\dfrac{1}{4}$ số vải đó. Hỏi xưởng còn lại bao nhiêu mét vải?
\shortans{3600}
\loigiai{
\textbf{Đáp số:} $3600$\\
May hết: $4800 : 4 = 1200\text{ m}$. Còn lại: $4800 - 1200 = 3600\text{ (m)}$.
}
\end{ex}

\begin{ex} %[3D2V58]
Tìm $x$, biết: $x : 5 + 1400 = 2600$.
\shortans{6000}
\loigiai{
\textbf{Đáp số:} $6000$\\
$x : 5 = 1200 \Rightarrow x = 6000$.
}
\end{ex}

\begin{ex} %[3D2V58]
Tổng của hai số là 6000, hiệu của hai số là 1000. Tìm số lớn.
\shortans{3500}
\loigiai{
\textbf{Đáp số:} $3500$\\
Số lớn là: $(6000 + 1000) : 2 = 3500$.
}
\end{ex}
